\section{Giới thiệu}
\label{sec:gioi-thieu}

\subsection{Bài toán và động cơ}

Thủ tục hộ tịch và chứng thực là nhóm thủ tục hành chính mà người dân tiếp xúc
nhiều nhất: khai sinh, khai tử, kết hôn, cấp trích lục, chứng thực bản sao và
chứng thực chữ ký. Thông tin về trình tự, thành phần hồ sơ, thời hạn giải quyết
và lệ phí đều được công bố công khai, nhưng phân tán trên nhiều văn bản ở nhiều
cấp hiệu lực khác nhau, và thay đổi theo thời gian. Người dân thường không biết
điều kiện áp dụng nào tương ứng với trường hợp của mình.

Một trợ lý hỏi đáp có thể giảm chi phí tra cứu, nhưng chỉ khi nó thỏa ba điều
kiện. Thứ nhất, câu trả lời phải nêu đúng điều kiện áp dụng, không chỉ nêu kết
luận chung. Thứ hai, câu trả lời phải \emph{dẫn nguồn} để người dùng tự kiểm tra,
vì giá trị pháp lý thuộc về văn bản gốc chứ không thuộc về mô hình. Thứ ba, khi
câu hỏi ngoài phạm vi hoặc chưa đủ căn cứ, trợ lý phải nói rõ thay vì đoán. Điều
kiện thứ ba quan trọng hơn điều thứ nhất về mặt rủi ro: một câu trả lời sai nhưng
tự tin gây hại nhiều hơn một lời từ chối.

Ba điều kiện này định hình toàn bộ thiết kế Phần~D. Chúng không đạt được bằng cách
huấn luyện mô hình cho đúng nhiều hơn, mà cần một tín hiệu học phân biệt được
``đúng và có dẫn nguồn'' với ``trôi chảy nhưng không kiểm được''. Đó chính là lý do
đề tài ghép SFT với tối ưu hóa theo ưu tiên: SFT dạy định dạng và nội dung, còn DPO
dạy thứ tự ưu tiên giữa hai câu trả lời mà con người đã xếp hạng.

\subsection{Phạm vi bắt buộc theo đề}

Đề tài 9 \cite[tr.~18--19]{de} yêu cầu bốn phần trên cùng một chủ đề, với trọng số
điểm như sau.

\begin{center}
\begin{tabular}{llr}
\toprule
\textbf{Phần} & \textbf{Nội dung} & \textbf{Điểm}\\
\midrule
A & Cài đặt từ đầu, có dẫn xuất và đối chiếu thư viện chuẩn & 20\\
B & Chứng minh chặt chẽ, đầy đủ giả thiết & 25\\
C & Bài tập tính toán, phân tích, thực nghiệm có kiểm soát & 15\\
D & Ứng dụng: thu thập/gán nhãn, tinh chỉnh, đánh giá & 30\\
\midrule
 & Báo cáo, tái lập, bảo vệ & 10\\
\bottomrule
\end{tabular}
\end{center}

Phần~D có trọng số lớn nhất. Đề cũng ràng buộc rằng ít nhất một cài đặt của
Phần~A phải được dùng trực tiếp hoặc so sánh trong quy trình Phần~D; Phần~D được
phép dùng thư viện chuẩn nhưng điều đó không thay thế nghĩa vụ tự cài đặt ở
Phần~A.

Hai phần lý thuyết B và C phải dựa trên năm bài báo trong danh mục tài liệu tham
khảo của đề: \cite{attention}, \cite{lora}, \cite{qlora}, \cite{rlhf},
\cite{dpo}. Báo cáo ghi rõ mục và phương trình được dùng ở từng chỗ.

\subsection{Đầu vào và đầu ra của sản phẩm}

Sản phẩm Phần~D là một trợ lý hỏi đáp.

\begin{description}[leftmargin=2em,style=nextline]
\item[Đầu vào] Một câu hỏi tiếng Việt về thủ tục hộ tịch hoặc chứng thực tại
Hà Nội, do người dùng tự diễn đạt. Không giả định người dùng biết tên thủ tục,
số hiệu văn bản hay thuật ngữ hành chính.
\item[Đầu ra mong muốn] Câu trả lời nêu đúng thông tin và điều kiện áp dụng, kèm
dẫn nguồn để người dùng kiểm tra. Nếu câu hỏi ngoài phạm vi hộ tịch -- chứng thực
tại Hà Nội, hoặc chưa đủ căn cứ trong dữ liệu đã thu thập, trợ lý nêu rõ điều đó
thay vì đoán.
\end{description}

Hai điều kiện ``dẫn nguồn'' và ``từ chối khi chưa đủ căn cứ'' là tiêu chí đánh giá
chính, không phải tính năng phụ. Chúng xuất hiện trong rubric chấm cặp ưu tiên
(Mục~\ref{sec:phan-d}) và trong giao thức đánh giá cuối.

\subsection{Cấu trúc báo cáo và quy ước trình bày}

Mục~\ref{sec:phan-a} trình bày Phần~A theo trình tự dẫn xuất $\to$ cài đặt $\to$
đối chiếu. Mục~\ref{sec:phan-b} là sáu chứng minh. Mục~\ref{sec:phan-c} là bốn
bài tập tính toán. Mục~\ref{sec:phan-d} là Phần~D, trình bày theo trạng thái thực
tế. Mục~\ref{sec:tai-lap} ghi điều kiện tái lập. Mục~\ref{sec:dao-duc} ghi các
ràng buộc đạo đức và pháp lý. Mục~\ref{sec:gioi-han} nêu giới hạn và hướng mở
rộng. Phụ lục~\ref{sec:phu-luc-ai} khai báo phạm vi dùng AI tạo sinh.

Ba quy ước áp dụng xuyên suốt:

\begin{enumerate}[leftmargin=2em]
\item \textbf{Mọi logarit là logarit tự nhiên}, kể cả trong perplexity và trong
các log-tỷ số của DPO.
\item \textbf{Nhãn trạng thái.} Mỗi hạng mục được gắn \dalam, \motphan{} hoặc
\chualam{} theo bằng chứng trong kho mã. Nhãn không phải dự báo; một hạng mục đã
lập kế hoạch nhưng chưa chạy vẫn là \chualam.
\item \textbf{Phân biệt số đo với số minh họa.} Những ví dụ số dùng để minh họa
phép tính được ghi rõ là minh họa. Chúng không được trình bày như dữ liệu đề chính
thức hay như kết quả thực nghiệm.
\end{enumerate}

\subsection{Phân công trong nhóm}

Nhóm không chia cứng bốn phần cho ba thành viên. Cao Văn Hà thực hiện bản đầu
xuyên suốt A--D, tổng hợp và nộp bài. Hà Mạnh Thế và Trần Duy Hưng luân phiên rà
soát từng đầu việc: chạy lại một phép kiểm mã, tự giải lại một bước chứng minh
hoặc tính toán, hoặc đối chiếu một bản ghi QA với URL và đúng mục trong nguồn.
Cách này nhằm bảo đảm cả ba thành viên đều giải thích được bất kỳ đoạn mã hoặc
chứng minh nào khi vấn đáp, theo yêu cầu phần~II của đề.

Riêng bước tạo nhãn ưu tiên cho Phần~D có quy trình khác và không luân phiên: Hà
Mạnh Thế và Trần Duy Hưng \emph{đều} chấm từng cặp một cách độc lập, không chia
đôi số cặp; Cao Văn Hà chỉ phân xử \emph{sau} khi nhãn độc lập đã bị khóa và mức
đồng thuận đã được đo. Lý do là yêu cầu của đề về hai người đánh giá độc lập chỉ
có nghĩa khi hai nhãn thực sự được tạo độc lập trên cùng một cặp.
